\documentclass{article}

\usepackage[spanish]{babel}
\usepackage{amsmath}
\usepackage[a4paper]{geometry}
\usepackage{parskip}
\usepackage{ragged2e}

\setlength{\parskip}{5mm}

\begin{document}

\begin{center}
 \textbf{\huge Examen De Física}   
\end{center}

\flushleft
\begin{enumerate}
    \item  considerem un objecte que té com a vector posició:
    
\begin{equation*}
\vec{r}=\left(\frac{1}{2}t^2+\frac{1}{3}t^3\right)
\end{equation*}

Trobeu els vectors velocitat i acceleració total i els seus mòduls. Trobeu també els mòduls de les components intrínseques de l'acceleració. Finalment, trobareu una expressió que relacioni el radii de curvatura amb el temps i a més, escriviu l'equació de la trajectòria
\item  Quina distància recorre en un periode complert la partícula que oscil·la amb amplitud A?
\item Si sabem que la velocitat d'un oscil·lador d'amplitud A és zero en certs instants, podem dir exactament quant val L'elongació en aquest instants
\item Quin és el valor d'un oscilador d'amplitud A i freqüència $f$ quan la seva velocitat és màxima? I quan el seu desplaçament és màxim?
\item Poden tenir el mateix sentit l'acceleració i el desplaçament d'un oscil·lador harmònic simple? I l'acceleració i la velocitat? I la velocitat i el desplaçament?
\item Una partícula té un desplaçament $x$ donat per:

\begin{equation*}
    x=0.3\cos{\left(2t+\frac{\pi}{6}\right)}
\end{equation*}

en unitats del sistema internacional. Trobeu:

\begin{enumerate}
     \item freqüència, periode i pulsació del moviment.
    \item on és la partícula per $t=1s$?
    \item velocitat i acceleració per un temps qualsevol $t$.
    \item velocitat, acceleració i posició inicials de la prtícula.
\end{enumerate}
\end{enumerate}

\end {document}